\subsection{Claim-level provenance and verification systems}

There is also direct prior art for representing and verifying claims at finer granularity than whole documents. Zhang, Ives and Roth formalize provenance graphs for natural-language claims in order to track where a claim came from and how it evolved \cite{zhang2020claim}. In biomedicine, the SEE framework represents scientific claims together with their provenance, argumentative relations, subjects, methods and assumptions in an RDF/OWL evidence model \cite{boelling2014see}. These systems make clear that claim-level provenance and evidence structure are not new functions introduced by Orion.

Verification systems add another layer. ProVe checks whether knowledge-graph triples are actually supported by the textual sources recorded as their provenance rather than assuming that a documented source entails the graph assertion \cite{amaral2024prove}. Open-domain scientific claim verification further shows that retrieval source and retrieval method materially affect evidence recall and precision when the relevant document is not already supplied to the verifier \cite{vladika2024}. In scholarly LLM interfaces, PaperTrail decomposes generated answers and sources into claim/evidence units to surface unsupported assertions and omissions \cite{martinboyle2026}.

Research-agent trajectories now supply an even closer operational parent. Zhuang et al. treat a completed experimental trajectory as candidate evidence rather than as knowledge, verify consequential artifacts before release, qualify cross-round intervention claims by execution validity and attribution, and retain only bounded Repairs or Guards while withholding unsupported findings \cite{zhuang2026trajectory}. Their industrial study also reports non-monotone adaptation---for 22 of 30 methods the last valid round was worse than an earlier best---and identifies affirmative applicability judgment as a remaining controller bottleneck. This owns substantial territory around verified proposal handoffs, negative-result retention, applicability-bounded experimental records and non-monotone research history. It also provides empirical motivation for preserving validated intermediate states rather than treating the final trajectory endpoint as authoritative.

These works narrow the present contribution in a useful way. Epistemic mechanics does not claim novelty for claim decomposition, provenance graphs, evidence attribution, source retrieval, support/refute classification, bounded artifact verification or applicability-qualified experimental records in isolation. Its additional object is the broader governed scientific-state transition around those operations: what kind of support was observed, under what context and independence assumptions, which scientific authority coordinates it can change, whether higher-order alignment obstructions and conflicting objects remain explicitly represented, and whether later discovery can reopen a state that had previously been judged saturated. Conversely, the present formal paper does not reproduce Zhuang et al.'s industrial intervention results or claim comparable empirical validation.
